% 让 Latex Workshop/TeX 引擎知道这是“主文件”

% 放在 \documentclass 之前
\PassOptionsToPackage{backend=biber,style=numeric-comp,sorting=nyt,refsegment=chapter,defernumbers=true}{biblatex}

\documentclass[lang=cn,scheme=chinese,12pt]{elegantbook}

% 显式加载 biblatex，确保使用 biber
\usepackage{biblatex}

%\usepackage[pass]{geometry}
\title{人工智能的数学思维}
%\subtitle{Elegant\LaTeX{} 经典之作}
%\author{吴敏 蒲戈光  }
\institute{华东师范大学 软件工程学院}
\date{2024}
%\version{4.3}
%\bioinfo{自定义}{信息}

%\extrainfo{不要以为抹消过去，重新来过，即可发生什么改变。--- 比企谷八幡}
% --- 全局：三层级（章→节→小节），编号与目录仅到小节 ---
\setcounter{secnumdepth}{2}
\setcounter{tocdepth}{2}
% 版式统一：将 \subsubsection 呈现为段落式标题（与此前章节内的局部设置一致）
\renewcommand\subsubsection{\paragraph}
%\logo{logo-blue.png}
\cover{cover.jpg}

%\setCJKmainfont{Songti SC}

% 本文档命令
\usepackage{array}
\newcommand{\ccr}[1]{\makecell{{\color{#1}\rule{1cm}{1cm}}}}
\usepackage{graphicx}
\usepackage{amsmath}  % 确保加载了 amsmath 包
\usepackage{amssymb}

\usepackage{titlesec} 
\usepackage{subcaption}
\usepackage{pgfplots}
\pgfplotsset{compat=1.18}

% --- 文本/URL 自动换行，减少 Overfull ---
\usepackage{xurl} % 让长链接可断行
\usepackage{microtype}           % 微排版（XeLaTeX 支持部分特性）
\sloppy                          % 宽松排版（先用，后面可再收紧）
% 临时缓解 Overfull \hbox（版面溢出）
\emergencystretch=4em

% --- 中文换行更友好 ---
\XeTeXlinebreaklocale "zh"
\XeTeXlinebreakskip = 0pt plus 0.2em minus 0.1em

% --- hyperref 选项务必在加载前传入 ---
\PassOptionsToPackage{unicode,breaklinks}{hyperref}
\usepackage{hyperref} % 可选：使得 \eqref 的引用带有链接功能
\hypersetup{hidelinks}

% --- 书签里禁用/替换某些命令（安全写法） ---
\usepackage{newunicodechar}
\newunicodechar{＝}{=}
\newunicodechar{－}{-}
\newunicodechar{Ｄ}{D}
\newunicodechar{ｇ}{g}
\newunicodechar{^^^^200b}{} % 忽略零宽空格 U+200B（注意用 4 个 ^ 的 16 进制写法）
% 推荐把罗马数字的 Unicode 映射掉（见第 4 节）：（避免 Missing character）
\newunicodechar{Ⅰ}{I}
\newunicodechar{Ⅱ}{II}
\newunicodechar{Ⅲ}{III}

% PDF 书签里不允许出现的东西（如 \\、数学），这里给替代
\pdfstringdefDisableCommands{%
  \def\eqref#1{(\ref{#1})}%
  \def\texttt#1{#1}% 书签里去等宽体
  \def\footnote#1{}% 书签里丢弃脚注
  % 如有别的自定义命令造成书签报警，再按此格式加一行
}
% --- 目录书签更稳，消除 “Rerun … use package ‘bookmark’” ---
\usepackage{bookmark}

\newcommand{\chapterwithsubtitle}[2]{ \chapter{#1\\ \textit{#2}}  }

% --- 中文右对齐名言（解决 cmr12 缺中文、避免 \\ 进书签）
\newcommand{\cnquote}[2]{%
  \begin{flushright}
    \begin{minipage}{0.85\linewidth}
      \kaishu #1\par\medskip
      \hfill---\,#2
    \end{minipage}
  \end{flushright}
}

\newenvironment{appendixsection}[1]{
    \section*{附录：#1}
    \addcontentsline{toc}{section}{附录：#1}
}{}


\newenvironment{story}[1]
{
    \section*{小故事：#1}
    \addcontentsline{toc}{section}{小故事：#1}
}{}

% 修改标题页的橙色带
% \definecolor{customcolor}{RGB}{32,178,170}
% \colorlet{coverlinecolor}{customcolor}

% 不需要 ifthen；IfFileExists 属于 LaTeX 内核
\makeatletter
\newcommand{\preferInput}[1]{%
  \begingroup
  \def\@base{#1}%
  \IfFileExists{\@base_polished.tex}{%
    \typeout{[preferInput] -> Using \@base_polished.tex}%
    \input{\@base_polished}%
  }{%
    \IfFileExists{\@base_refined.tex}{%
      \typeout{[preferInput] -> Using \@base_refined.tex}%
      \input{\@base_refined}%
    }{%
      \typeout{[preferInput] -> Fallback to \@base.tex}%
      \input{\@base}%
    }%
  }%
  \endgroup
}
\makeatother



\begin{document}
\maketitle

\frontmatter
% 前言/致谢/目录等（无编号）

\preferInput{chapters/00自序}
\preferInput{chapters/01他序}

\tableofcontents

\mainmatter   % ★ 从这里开始：第1章、第2章……

\preferInput{chapters/0-Intro.v2}

% ===== 正文章节 =====
\part{追问智能}
\preferInput{chapters/1-Def.v1}  % 什么是智能
\preferInput{chapters/2-History-Math.v3.2-conservative} % 智能，生于数学
% \preferInput{chapters/2-History-Math.v3.1} % 智能，生于数学（旧版本，备选）
\preferInput{chapters/3-History-computer.v1}  % **MOD** 更新至最新存在版本（v1）

\part{学习的逻辑}
\preferInput{chapters/3-Thinking.v1}  % **MOD** 更新版本 - 三种思维方法论
\preferInput{chapters/4-Tools.v1}  % **MOD** 更新版本 - 数学工具基础
\preferInput{chapters/5-ML.v3}  
\preferInput{chapters/6ML2DL.v1} % **MOD** 更新版本 - 从机器学习到深度学习
\preferInput{chapters/7-DL.v1} % **MOD** 更新版本 - 深度学习实践
%\chapter*{人工智能的其他方向及推荐阅读}

本书聚焦于人工智能的数学与思维逻辑，尚未覆盖以下内容：
\begin{itemize}
\item {\bf  Transformer 模型与语言智能：}推荐《Attention is All You Need》

\item{强化学习与决策系统：}推荐 Sutton $\&$ Barto 的《Reinforcement Learning》

\item{图神经网络与结构化学习：}
\end{itemize}


读者若对这些方向感兴趣，可由此继续探索。

\part{理解的边界}
\preferInput{chapters/8-深度学习的挑战与幻觉.v4} % **MOD** 二次扩写更新版本
\preferInput{chapters/9-AI的边界：图灵机、哥德尔与计算的极限.v3} % **MOD** 二次扩写更新版本
\preferInput{chapters/10-符号、神经与因果：智能的三种建构方式.v3} % **MOD** 二次扩写更新版本
\preferInput{chapters/11-智能的物理基础：从熵到量子计算.v3} % **MOD** 二次扩写更新版本

\part{回归与超越}
\preferInput{chapters/12-语言中的智能：Transformer的意外胜利.v3} 
% **MOD** 二次扩写更新版本
\preferInput{chapters/13-AI+.v3} 
% **MOD** 二次扩写更新版本
\preferInput{chapters/14-理解智能的尽头：理性与取舍的未来.v3} % **MOD** 二次扩写更新版本
\preferInput{chapters/15-伦理与挑战.v3} % **MOD** 二次扩写更新版本

\part{向上的道路}
\preferInput{chapters/16-Conclusion.v1} % **MOD** 二次扩写更新版本

\nocite{*}
\printbibliography[heading=bibintoc, title=参考文献]

\end{document}
